\chapter*{Introduction Générale}
\addcontentsline{toc}{chapter}{Introduction Générale}

\section*{Contexte socio-économique et technologique}
L'artisanat marocain, riche de plusieurs siècles d'histoire, constitue l'un des piliers fondamentaux de l'économie nationale et du patrimoine culturel immatériel. Vecteur d'identité et de savoir-faire transmis de génération en génération, ce secteur fait cependant face, à l'ère de la mondialisation, à des défis structurels majeurs. La transition numérique, qui a bouleversé les habitudes de consommation à l'échelle globale, impose aujourd'hui une redéfinition des canaux de distribution traditionnels. 

Bien que les produits artisanaux (qu'ils soient issus de la poterie, de la maroquinerie, du travail du bois ou de la dinanderie) jouissent d'une réputation d'excellence à travers le monde, l'accès direct aux marchés reste souvent entravé. Les artisans, principalement regroupés en coopératives ou opérant de manière individuelle, dépendent fréquemment d'une longue chaîne d'intermédiaires, réduisant ainsi leur marge bénéficiaire et limitant leur visibilité auprès du consommateur final.

C'est dans ce contexte de fracture numérique qu'émerge la nécessité impérieuse de concevoir des solutions logicielles adaptées, capables de démocratiser l'accès au commerce électronique pour les acteurs de l'artisanat.

\section*{Problématique et Objectifs du Projet}
La problématique centrale de notre projet de fin d'études s'articule autour de la question suivante : \textit{Comment tirer parti des technologies du web moderne pour concevoir une plateforme de mise en relation directe, sécurisée et ergonomique, permettant la valorisation et la commercialisation des produits artisanaux sans la barrière des intermédiaires ?}

Pour répondre à cette problématique, le projet \textbf{Lartisan} a été initié. Il s'agit d'une plateforme d'intermédiation numérique (Marketplace) conçue spécifiquement pour l'écosystème artisanal. Les objectifs de ce projet sont multidimensionnels :
\begin{itemize}[label=\textbullet]
    \item \textbf{Objectif économique} : Maximiser le profit direct de l'artisan en instaurant un modèle \textit{Peer-to-Peer} (Vendeur-Acheteur).
    \item \textbf{Objectif technologique} : Déployer une architecture logicielle robuste, scalable et hautement sécurisée en exploitant les frameworks de dernière génération.
    \item \textbf{Objectif social} : Promouvoir le patrimoine culturel marocain à travers une interface visuelle immersive et moderne, respectant les standards actuels de l'expérience utilisateur (UX).
\end{itemize}

\section*{Démarche Méthodologique et Structure du Rapport}
Afin de mener à bien ce projet d'ingénierie logicielle, nous avons adopté une démarche itérative hybride. Cette approche s'appuie sur la rigueur de la modélisation orientée objet pour la structuration des données, combinée à une analyse processus détaillée pour la définition des règles métier.

Le présent document synthétise l'ensemble de notre travail et s'organise en deux chapitres principaux :
\begin{enumerate}[label=\arabic*.]
    \item \textbf{Le Chapitre 1 (Méthode de travail)} est consacré à l'étude préalable, à la définition exhaustive du cahier des charges et à la modélisation conceptuelle du système à l'aide des standards UML et MCT.
    \item \textbf{Le Chapitre 2 (Application Finale)} aborde les choix technologiques de réalisation, justifie l'architecture technique retenue, et présente en détail les interfaces graphiques constituant la solution finale livrable.
\end{enumerate}
Une conclusion générale viendra clôturer ce rapport en dressant le bilan des compétences acquises et en ouvrant sur les perspectives d'évolution futures du projet Lartisan.
